% Generated draft: requires mathematical and source review.
\documentclass[11pt]{article}
\usepackage[margin=25mm]{geometry}
\usepackage{fontspec}
\setmainfont{Latin Modern Roman}
\usepackage{amsmath,amssymb,mathtools,unicode-math}
\setmathfont{Latin Modern Math}
\usepackage{enumitem,needspace,xurl,hyperref}
\setlength{\parindent}{0pt}
\setlength{\parskip}{5pt}
\setlength{\emergencystretch}{3em}
\newcommand{\N}{\mathbb{N}}
\newcommand{\Z}{\mathbb{Z}}
\newcommand{\Q}{\mathbb{Q}}
\newcommand{\R}{\mathbb{R}}
\newcommand{\C}{\mathbb{C}}
\newcommand{\F}{\mathbb{F}}
\newcommand{\A}{\mathbb{A}}
\newcommand{\PP}{\mathbb{P}}
\newcommand{\E}{\mathbb{E}}
\newcommand{\eps}{\varepsilon}
\newcommand{\dd}{\,\mathrm{d}}
\begin{document}
\section*{UnsolvedMath ID: 30001667}
\textbf{Source problem code:} catalogue entry
\subsection{Definitions and mathematical statement}
For a positive integer \(d\) and a set \(V\), a \(d\)-uniform set system is a family \(\mathcal{S}\subseteq\binom{V}{d}\), where \(\binom{V}{d}\) denotes the family of all \(d\)-element subsets of \(V\). This terminology alone does not specify the cofilling profile in the supplied record.

The record describes \(\phi_d(\alpha)\) as a cofilling profile defined using minimal \(d\)-uniform systems and their coboundaries. It asks whether the asserted upper bound
\[
\phi_d(\alpha)\leq \frac{d}{d+1}\alpha
\]
is optimal, and whether the corresponding extremal example is essentially unique.

The supplied record does not define the coboundary operator, minimality, or \(\phi_d\). It also omits the admissible values of \(d\) and \(\alpha\), density normalizations, finite or asymptotic conventions, the extremal construction, and the equivalence intended by essential uniqueness. Consequently, the displayed question cannot be made into a fully quantified conjecture from this record alone. In particular, optimality is not silently interpreted as pointwise equality, and essential uniqueness is not assigned an equivalence relation.
\subsection{Short English statement}
The supplied question asks whether the stated linear upper bound for a cofilling profile is sharp and whether its extremal construction is essentially unique. The profile, parameter ranges, construction, and uniqueness criterion are missing, so a precise formulation requires additional source information. Supplied dated curation labels the problem open, including for \(d=2\); this status has not been independently verified.
\subsection{Sources}
\begin{itemize}
\item Supplied local UnsolvedMath registry record: permanent ID \(30001667\), title \emph{Optimal Cofilling Profiles for Uniform Set Systems}, exact source problem code \texttt{catalogue entry}, in the Oberwolfach Reports Open Problems collection. The record gives proposed year 2011 and names no individual proposer.
\item \emph{Combinatorics} (2011), Oberwolfach Reports, pp.~5--83, DOI 10.4171/owr/2011/01: \url{https://doi.org/10.4171/owr/2011/01}. The supplied background describes these reports as records of mathematical workshops and their open-problem discussions. No pinpoint passage or complete defining text is supplied here.
\item Supplied literature reference: Matou\v{s}ek--Wagner, \url{https://arxiv.org/abs/1102.3515}. The record supplies no title or pinpoint location for this reference.
\item Literature assessment dated exactly \texttt{2026-08-21}: the supplied curation states that the primary report leaves the optimality and uniqueness conjecture open, even for \(d=2\), and lists the two references above as checked. This is supplied dated curation, not an independently verified literature assessment; the record expressly requests independent verification of time-sensitive status and cited claims.
\item The registry labels the statement \texttt{corrected\_verified} and reports that the official report was checked and the statement repaired to match its source passage. These are supplied provenance claims, not verification performed for this draft.
\end{itemize}
\end{document}
